\documentclass[../../main]{subfiles}

\begin{document}

\section{Espaços de Funções}
\label{sec:matematica:funcoes:espacos-de-funcoes}

\subsection{Definição}

\begin{defn}[Espaço de Funções]
	\label{def:matematica:funcoes:espacos:espaco-funcoes}

	Sejam \(A\) e \(B\) conjuntos.

	O espaço de funções de \(A\) em \(B\) é o conjunto

	\[
		B^A
		=
		\{
		f
		\mid
		f:A\to B
		\}.
	\]

	Ou seja, \(B^A\) é o conjunto de todas as funções com domínio
	\(A\) e contradomínio \(B\).

\end{defn}

\begin{obs}

	Os elementos de \(B^A\) são funções.

\end{obs}

\begin{ex}

	Seja

	\[
		A=\{1,2\}
	\]

	e

	\[
		B=\{0,1\}.
	\]

	Então

	\[
		B^A
	\]

	possui quatro elementos:

	\[
		\begin{aligned}
			f_1(1) & =0, & f_1(2) & =0, \\
			f_2(1) & =0, & f_2(2) & =1, \\
			f_3(1) & =1, & f_3(2) & =0, \\
			f_4(1) & =1, & f_4(2) & =1.
		\end{aligned}
	\]

\end{ex}

\subsection{Notação Exponencial}

\begin{obs}

	A notação

	\[
		B^A
	\]

	não é arbitrária.

	Quando \(A\) e \(B\) são conjuntos finitos,

	\[
		|B^A|
		=
		|B|^{|A|}.
	\]

\end{obs}

\begin{prop}
	\label{prop:matematica:funcoes:espacos:cardinalidade-finita}

	Se \(A\) e \(B\) são conjuntos finitos, então

	\[
		|B^A|
		=
		|B|^{|A|}.
	\]

\end{prop}

\begin{proof}

	Para cada elemento de \(A\) existem

	\[
		|B|
	\]

	possíveis escolhas para sua imagem.

	Como as escolhas são independentes,

	pelo princípio multiplicativo,

	\[
		|B^A|
		=
		|B|^{|A|}.
	\]

\end{proof}

\begin{ex}

	Como

	\[
		|\{0,1\}^{\{1,2,3\}}|
		=
		2^3,
	\]

	existem exatamente

	\[
		8
	\]

	funções de

	\[
		\{1,2,3\}
	\]

	em

	\[
		\{0,1\}.
	\]

\end{ex}

\subsection{Função de Avaliação}

\begin{defn}[Função de Avaliação]
	\label{def:matematica:funcoes:espacos:avaliacao}

	Sejam \(A\) e \(B\) conjuntos.

	Define-se a função de avaliação

	\[
		\operatorname{ev}
		:
		B^A\times A
		\to
		B
	\]

	por

	\[
		\operatorname{ev}(f,x)
		=
		f(x).
	\]

\end{defn}

\begin{prop}

	A função de avaliação está bem definida.

\end{prop}

\begin{proof}

	Se

	\[
		(f,x)
		\in
		B^A\times A,
	\]

	então

	\[
		f:A\to B.
	\]

	Logo

	\[
		f(x)\in B.
	\]

	Portanto

	\[
		\operatorname{ev}(f,x)
	\]

	é um elemento bem definido de \(B\).

\end{proof}

\begin{obs}

	A avaliação é a operação fundamental que permite recuperar valores
	de uma função considerada como elemento do espaço \(B^A\).

\end{obs}

\subsection{Caracterização Extensional}

\begin{prop}[Extensionalidade]
	\label{prop:matematica:funcoes:espacos:extensionalidade}

	Se

	\[
		f,g\in B^A,
	\]

	então

	\[
		f=g
	\]

	se, e somente se,

	\[
		(\forall x\in A)
		\;
		f(x)=g(x).
	\]

\end{prop}

\begin{proof}

	Trata-se do princípio da extensionalidade para funções.

\end{proof}

\subsection{Espaços de Funções Constantes}

\begin{defn}[Função Constante]
	\label{def:matematica:funcoes:espacos:constante}

	Seja

	\[
		b\in B.
	\]

	A função constante associada a \(b\) é a função

	\[
		c_b:A\to B
	\]

	definida por

	\[
		c_b(x)=b
	\]

	para todo

	\[
		x\in A.
	\]

\end{defn}

\begin{obs}

	Cada elemento de \(B\) determina uma função constante em \(B^A\).

\end{obs}

\end{document}
